\documentclass[aspectratio=169,10pt]{beamer}
\usepackage{plan4ari_beamer}

\usepackage{iftex}
\ifPDFTeX%
  \usepackage[T1]{fontenc}
  \usepackage[utf8]{inputenc}
  % Times-compatible fallback for pdfLaTeX.
  \usepackage{newtxtext,newtxmath}
\else
  \usepackage{fontspec}
  % Prefer one of the reference fonts named in the official template.
  \IfFontExistsTF{Times New Roman}{%
    \setmainfont{Times New Roman}
  }{%
    \IfFontExistsTF{Nimbus Roman No9 L}{%
      \setmainfont{Nimbus Roman No9 L}
    }{%
      \IfFontExistsTF{Nimbus Roman}{%
        \setmainfont{Nimbus Roman}
      }{%
        % Metric-compatible Times-style fallback; verify against the official template before submission.
        \setmainfont{TeX Gyre Termes}
      }%
    }%
  }
\fi
\usepackage{lmodern}
\usepackage{graphicx}
\usepackage{microtype}
\usepackage{ragged2e}

\usepackage{amsmath,amssymb}
\usepackage{tikz}
\usetikzlibrary{arrows,positioning,shapes.geometric,fit,calc,matrix,shadows.blur}
\usepackage{pgfgantt}
\PassOptionsToPackage{table}{xcolor}
\usepackage{xcolor}
\hypersetup{colorlinks=true,linkcolor=planblue,urlcolor=planblue,citecolor=planblue}

\usepackage[
    backend=biber,
    style=verbose,
    autocite=footnote,
    doi=true,
    url=false
]{biblatex}


\usepackage[marginal]{footmisc}


% ---------- Tables, figures, lists ----------
\usepackage{booktabs,tabularx,array,multirow}
\usepackage{longtable}
\usepackage{makecell}
\usepackage{subcaption}
\usepackage{caption}
\captionsetup{labelfont=bf}


% ---------- Header/footer and page-count warning ----------
\usepackage{empheq}
\usepackage[most]{tcolorbox}
\tcbuselibrary{breakable}

\usepackage[normalem]{ulem}

\usepackage{xparse}

\usepackage{xltabular}
\keepXColumns

% Compatibility for longtable v4.24 and older: its output routine used
% infinitely shrinkable \vss glue, which LuaTeX reports when a long table splits.
\makeatletter
\@ifpackagelater{longtable}{2026/05/31}{}{%
  \ExplSyntaxOn
  \def\LT@output{%
    \ifnum\outputpenalty <-\@Mi
      \ifnum\outputpenalty > -\LT@end@pen
        \LT@err{floats~ and~ marginpars~ not~ allowed~ in~ a~ longtable}\@ehc
      \else
        \setbox\z@\vbox{\unvbox\@cclv}%
        \ifdim \ht\LT@lastfoot>\ht\LT@foot
          \dimen@\pagegoal
          \advance\dimen@\ht\LT@foot
          \advance\dimen@-\ht\LT@lastfoot
          \ifdim\dimen@<\ht\z@
            \setbox\@cclv\vbox{%
              \unvbox\z@\copy\LT@foot
              \vskip 0pt plus \maxdimen minus \normalbaselineskip
            }%
            \@makecol
            \@expl@@@mark@update@singlecol@structures@@
            \@outputpage
            \global\vsize\@colroom
            \setbox\z@\vbox{\box\LT@head}%
          \fi
        \fi
        \unvbox\z@\box\ifvoid\LT@lastfoot\LT@foot\else\LT@lastfoot\fi
        \UseTaggingSocket{tbl/longtable/foot}%
      \fi
    \else
      \setbox\@cclv\vbox{%
        \unvbox\@cclv\copy\LT@foot
        \vskip 0pt plus \maxdimen minus \normalbaselineskip
      }%
      \UseTaggingSocket{tbl/longtable/foot}%
      \@makecol
      \@expl@@@mark@update@singlecol@structures@@
      \@outputpage
      \global\vsize\@colroom
      \copy\LT@head\nobreak
    \fi}%
  \ExplSyntaxOff
}
\makeatother

\usepackage{mdframed}
\usepackage{wrapfig}


% Keep normal, uncompressed character spacing and standard single spacing.
\linespread{1.0}\selectfont
\setlength{\parindent}{0pt}
\setlength{\parskip}{0.35em}
\raggedbottom%

\setlength{\tabcolsep}{2.2pt}


% pa glossary colour
\definecolor{paWhite}{HTML}{FFFFFF}

\definecolor{ink}{HTML}{16324A}
\definecolor{accent}{HTML}{087E8B}

\definecolor{paBlue}{HTML}{244A73}
\definecolor{paMidBlue}{HTML}{5287C0}
\definecolor{paLightBlue}{HTML}{EEF4F8}
\definecolor{paTermBlue}{HTML}{DCEAF3}
\definecolor{paBorder}{HTML}{9EB6C8}

\definecolor{paGreen}{HTML}{438961}
\definecolor{paMidGreen}{HTML}{5A966F}
\definecolor{paTermGreen}{HTML}{DCEFE3}
\definecolor{paLightGreen}{HTML}{EEF7F1}
\definecolor{paGreenBorder}{HTML}{A4C5B1}

\definecolor{paGray900}{HTML}{222222}
\definecolor{paGray700}{HTML}{4D5661}
\definecolor{paGray500}{HTML}{7B8490}
\definecolor{paGray300}{HTML}{C9D0D6}
\definecolor{paGray100}{HTML}{F3F5F7}

\definecolor{paOrange}{HTML}{D97706} 
\definecolor{paLightOrange}{HTML}{FDF4E7} 
\definecolor{paTermOrange}{HTML}{F6E2BF} 
\definecolor{paOrangeBorder}{HTML}{D8B070}

\definecolor{LightOrange}{HTML}{FFB347}
\definecolor{Orange}{HTML}{FFA500}
\definecolor{DarkOrange}{HTML}{FF8C00}
\definecolor{BurntOrange}{HTML}{CC5500}
\definecolor{paRed}{HTML}{C0392B}


% Keep the shared plan4ari_beamer colours, title sizes, margins and footer.
\setlength{\tabcolsep}{4pt}
\renewcommand{\arraystretch}{1.15}
\emergencystretch=1em

% ---------- Useful column types ----------
\newcolumntype{Y}{>{\raggedright\arraybackslash}X}
\newcolumntype{C}[1]{>{\centering\arraybackslash}p{#1}}
\newcolumntype{L}[1]{>{\raggedright\arraybackslash}p{#1}}
\newcolumntype{M}[1]{>{\centering\arraybackslash\bfseries}m{#1}}
\newcolumntype{J}[1]{>{\arraybackslash}p{#1}}

\newcolumntype{s}{>{\hsize=.8\hsize\raggedright\arraybackslash}X}

% ---------- pa concept box ----------
\newtcolorbox{paConcept}[1]{
    colback=paLightBlue,
    colframe=paBlue,
    colbacktitle=paBlue,
    coltitle=white,
    fonttitle=\bfseries,
    title={#1},
    boxrule=0.5pt,
    arc=1mm,
    left=1.5mm,
    right=1.5mm,
    top=1.5mm,
    bottom=1.5mm,
    before skip=3pt,
    after skip=3pt,
    breakable,
    enhanced
}

% ---------- pa concept box ----------
\newtcolorbox{paBox}[1][]{
    colback=paLightBlue,
    colframe=paBlue,
    boxrule=0.5pt,
    arc=1mm,
    left=1.5mm,
    right=1.5mm,
    top=1.5mm,
    bottom=1.5mm,
    before skip=3pt,
    after skip=3pt,
    breakable,
    enhanced,
    #1
}



% ---------- pa concept box ----------
\newtcolorbox{paBoxGreen}[1][]{
    colback=paLightGreen,
    colframe=paGreen,
    boxrule=0.1pt,
    arc=0.5mm,
    left=1.5mm,
    right=1.5mm,
    top=1.5mm,
    bottom=1.5mm,
    before skip=6pt,
    after skip=6pt,
    breakable,
    enhanced,
    #1
}

% ---------- pa examnple box ----------
\newtcolorbox{paExampleBox}[1][]{
    fonttitle=\bfseries,
    fontupper=\itshape,
    colback=paLightBlue,
    colframe=paLightBlue,
    coltitle=paBlue,
    boxrule=0.1pt,
    arc=0.1mm,
    left=1.mm,
    right=1.1mm,
    top=1.mm,
    bottom=1.mm,
    before skip=3pt,
    after skip=3pt,
    breakable,
    enhanced,
    #1
}


% ---------- pa light box ----------
\newtcolorbox{paLightBox}[1][]{
    colback=paGray100,
    colframe=paGray100,
    boxrule=0.1pt,
    arc=0.1mm,
    left=1.mm,
    right=1.1mm,
    top=1.mm,
    bottom=1.mm,
    before skip=3pt,
    after skip=3pt,
    breakable,
    enhanced,
    #1
}

% ---------- pa light box ----------
\newtcolorbox{paLightBluesBox}[1][]{
    colback=paLightBlue,
    colframe=paLightBlue,
    boxrule=0.1pt,
    arc=0.1mm,
    left=1.mm,
    right=1.1mm,
    top=1.mm,
    bottom=1.mm,
    before skip=3pt,
    after skip=3pt,
    breakable,
    enhanced,
    #1
}


\newtcbox{\boxedmath}[1][]{%
    nobeforeafter, math upper, tcbox raise base,
    enhanced, boxrule=0.2pt,
    #1}


\renewcommand{\cellalign}{vh}
\renewcommand{\theadalign}{vh}

\input{styles/math.tex}
% chktex-file 44



%%%%%%%%%%%%%%%%%%%%%%%%%%%%%%%%%%%%%%%%%%%%%%%%%%%%%%%%%%%%
% Basic Style Macros
\newcommand{\paemph}[1]{%
    \textcolor{pagRed}{\textit{#1}}%
}

\newcommand{\patextbf}[1]{%
    \textcolor{pagBlue}{\textbf{#1}}%
}

\newcommand{\patextbfo}[1]{%
    \textcolor{pagOrange}{\textbf{#1}}%
}
\newcommand{\placeholder}[1]{\textit{[#1]}}

\newcommand{\blueitem}{%
\textcolor{pagBlue}{\rule{0.8ex}{0.8ex}}% 
}
%%%%%%%%%%%%%%%%%%%%%%%%%%%%%%%%%%%%%%%%%%%%%%%%%%%%%%%%%%%%



%%%%%%%%%%%%%%%%%%%%%%%%%%%%%%%%%%%%%%%%%%%%%%%%%%%%%%%%%%%%
% Table Formatting
\newcommand{\parowcolored}[1]{%
    \rowcolor{#1}
    [\dimexpr\tabcolsep+0.2pt\relax]
    [\dimexpr\tabcolsep+0.2pt\relax]
}

\newcommand{\parowcolor}{%
    \parowcolored{pagMidBlue}%
}

\newcommand{\paheaderstrut}{%
  \rule[\dimexpr-\dp\strutbox-2pt\relax]{0pt}%
       {\dimexpr\ht\strutbox+\dp\strutbox+4pt\relax}%
}

\NewExpandableDocumentCommand{\pahcell}{o O{|l|} m}{%
    \IfNoValueTF{#1}{%
        \leavevmode\color{white}\paheaderstrut\textbf{#3}
    }{%
        \multicolumn{#1}{#2}{%
            \leavevmode\color{white}\paheaderstrut\textbf{#3}%
        }%
    }%
}

\NewDocumentCommand{\patableformat}{O{1.0} O{2pt} O{2pt}}{%
    \small
    \centering
    \renewcommand{\arraystretch}{#1}
    \setlength{\tabcolsep}{#2}
    \setlength{\extrarowheight}{#3}
    \arrayrulecolor{pagBorder}
}


\NewDocumentCommand{\patableformatnormalsize}{O{1.0} O{2pt} O{2pt}}{%
    \normalsize
    \centering
    \renewcommand{\arraystretch}{#1}
    \setlength{\tabcolsep}{#2}
    \setlength{\extrarowheight}{#3}
    \arrayrulecolor{pagBorder}
}

\newcommand{\rowspace}{\rule{0pt}{2ex}}


\NewDocumentCommand{\pafirstlinecell}{O{t} m m}{%
    \multirow[#1]{#2}{=}{%
        \parbox[t]{\dimexpr\linewidth-2\tabcolsep\relax}{%
            \raggedright #3 
        }%
    }% 
}

%%%%%%%%%%%%%%%%%%%%%%%%%%%%%%%%%%%%%%%%%%%%%%%%%%%%%%%%%%%%



%%%%%%%%%%%%%%%%%%%%%%%%%%%%%%%%%%%%%%%%%%%%%%%%%%%%%%%%%%%%
% TOC Formatting
\newlength{\palabelwidth}

\NewDocumentCommand{\pasubsubsectiononeline}{O{pagMidGreen} m +g o}{%
  \par
  \begingroup
  \raggedright
  \settowidth{\palabelwidth}{\textbf{#2}}%
  \addtolength{\palabelwidth}{4pt}%
  \phantomsection
  \addcontentsline{toc}{subsubsection}{#2}%
  \nopagebreak[4]%
  \noindent
  \IfNoValueTF{#3}{%
    \begin{tabularx}{\linewidth}{%
      >{\raggedright\arraybackslash}X}%
      \parowcolored{#1}\pahcell{#2}\\%
    \end{tabularx}%
  }{%
    \begin{tabularx}{\linewidth}{%
      >{\raggedright\arraybackslash}p{\palabelwidth}
      >{\raggedright\arraybackslash}X}%
      \cellcolor{#1}\pahcell{#2}
      &%
      \leavevmode\color{#1}\paheaderstrut\textbf{#3}
      \IfValueT{#4}{%
        \hfill
        {\color{black}#4} 
        \kern4pt \vrule width 1.5pt
      }\\%
      \IfNoValueT{#4}{\arrayrulecolor{#1}\hline}
    \end{tabularx}%
    \arrayrulecolor{pagBorder}%
  }%
  \par
  \endgroup
}

\NewDocumentCommand{\pasubsubsectionboxed}{O{pagMidGreen} m +g}{%
  \par
  \begingroup
  \phantomsection
  \addcontentsline{toc}{subsubsection}{#2}%
  \nopagebreak[4]%
  \noindent
  \begin{tabular}{|>{\raggedright\arraybackslash}%
    p{\dimexpr\linewidth-2\tabcolsep-2\arrayrulewidth\relax}|}%
    \rowcolor{#1}%
    \leavevmode\color{white}\paheaderstrut\textbf{#2}\\%
    \IfNoValueF{#3}{%
      \leavevmode\color{#1}\paheaderstrut\textbf{#3}\\%
      \hline
    }%
  \end{tabular}%
  \arrayrulecolor{pagBorder}%
  \par
  \endgroup
}

\NewDocumentCommand{\pasubsubsectionnamedboxed}
  {O{pagMidBlue} O{pagMidBlue} m m m o}{%
    \par
    \begingroup
    \raggedright
    \settowidth{\palabelwidth}{\textbf{#3}}%
    \addtolength{\palabelwidth}{4pt}%
    \IfNoValueT{#6}{
        \phantomsection
        \addcontentsline{toc}{subsubsection}{#3: #4}%
        \nopagebreak[4]%
    }
    \noindent
    \arrayrulecolor{#1}
    \begin{tabularx}{\dimexpr\linewidth-2pt\relax}{%
        >{\raggedright\arraybackslash}p{\palabelwidth}
        >{\raggedright\arraybackslash}X}%
        \cellcolor{#1}\pahcell{#3} &
        \leavevmode\color{#1}\paheaderstrut\textbf{#4}%
        \par
        \hrule height \arrayrulewidth
        \kern1pt
        \color{#2}\noindent\paheaderstrut#5\\%
    \end{tabularx}%
    \arrayrulecolor{pagBorder}%
    \par
    \endgroup
}
\NewDocumentCommand{\pasubsubsection}{m o}{%
    \par
    \phantomsection
    \addcontentsline{toc}{subsubsection}{#1}%
    \nopagebreak[4]%
    \textcolor{pagBlue}{%
        \textbf{#1}%
        \IfValueT{#2}{: \textbf{#2}}%
        \IfNoValueT{#2}{.}% 
    }% 
}

\newcounter{paRQ}
\NewExpandableDocumentCommand{\paRQrow}{m m}{%
  \IfValueT{#2}{%
    \IfBlankF{#2}{%
      \stepcounter{paRQ}%
      \patextbf{RQ\arabic{paRQ}} &
      \leavevmode\textit{#2}\\%
    }%
  }%
}

\newlength{\padefaultlabelwidth}
\NewDocumentCommand{\pasubsubsectionobjective}
  {O{pagMidGreen} m m m o o o o o}{%
    \par
    \begingroup
    \raggedright
    \settowidth{\palabelwidth}{\textbf{#2}}%
    \settowidth{\padefaultlabelwidth}{\textbf{RQ10}}%
    \ifdim\palabelwidth<\padefaultlabelwidth%
        \setlength{\palabelwidth}{\padefaultlabelwidth}%
    \fi%
    \addtolength{\palabelwidth}{4pt}%
    \phantomsection
    \addcontentsline{toc}{subsubsection}{#2: #3}%
    \nopagebreak[4]%
    \noindent
    \arrayrulecolor{#1}%
    \setlength{\extrarowheight}{1pt}%
    \begin{tabularx}{\linewidth}{%
        >{\raggedright\arraybackslash}p{\palabelwidth}
        >{\raggedright\arraybackslash}X}%
    \cellcolor{#1}\pahcell{#2} & 
    \leavevmode\color{#1}\paheaderstrut\textbf{#3}%
        \par
        \hrule height \arrayrulewidth
        \kern1pt
        \noindent\paheaderstrut#4\\%
        \paRQrow{#1}{#5}%
        \paRQrow{#1}{#6}%
        \paRQrow{#1}{#7}%
        \paRQrow{#1}{#8}%
        \paRQrow{#1}{#9}%
    \end{tabularx}%
    \arrayrulecolor{pagBorder}%
    \par
    \endgroup
}

%%%%%%%%%%%%%%%%%%%%%%%%%%%%%%%%%%%%%%%%%%%%%%%%%%%%%%%%%%%%%%%%

%%%%%%%%%%%%%%%%%%%%%%%%%%%%%%%%%%%%%%%%%%%%%%%%%%%%%%%%%%%%




%%%%%%%%%%%%%%%%%%%%%%%%%%%%%%%%%%%%%%%%%%%%%%%%%%%%%%%%%%%%
% ---------- Project ID ----------
\NewDocumentCommand{\patextbfgraphic}{O{pagGreen} O{pagGreen} m m}{%
\IfBlankTF{#3}{%
  \begin{tikzpicture}[baseline=(txt.base)]
    \node[inner xsep=0pt, inner ysep=0pt, text=#2, font=\bfseries] (txt) {#4};
    \draw[#1, line width=0.8pt] (txt.south west) -- (txt.south east);
  \end{tikzpicture}%
}{%
  \begin{tikzpicture}[baseline=(txt.base)]
    \node[inner sep=2pt, fill=#1, text=white, font=\bfseries] (tag) {#3};
    \node[anchor=west, inner xsep=0pt, inner ysep=0pt, text=#2, font=\bfseries] (txt) at ([yshift=-0.4pt]tag.east) {#4};
    \draw[#1, line width=0.8pt] (tag.south west) -- (txt.south east);
  \end{tikzpicture}%
}%
}

\newcommand{\ProjectName}{SMARTHER}
\newcommand{\CallIdentifier}{[HORIZON-EIC-2026-PATHFINDERCHALLENGES-01--03]}
\newcommand{\CallName}{[DeepRAP:\@ Deep Reasoning, Abstraction \& Planning]}
% \newcommand{\ProposalTitle}{%
% \begin{center}
% \begin{minipage}{0.15\textwidth}
% \includegraphics[width=\textwidth]{images/logo.png}
% \end{minipage}
% \begin{minipage}{0.7\textwidth}
% \centering
%   %\makebox[\textwidth][c]{%
%     \patextbfgraphic[pagOrange][black]
%       {}
%       {\patextbfo{S}emantic \patextbfo{M}etacognitive \patextbfo{A}utonomy Ha\patextbfo{R}ness for \patextbfo{T}rustworthy,}%
%   %}
%   \\[6pt]
%   %\makebox[\textwidth][c]{%
%     \patextbfgraphic[pagOrange][black]
%       {}
%       {\patextbfo{H}uman-Aware, \patextbfo{E}xplainable \patextbfo{R}obotics}%
%   %}%
% \end{minipage}
% \end{center}
% }

\newcommand{\ProposalTitle}{%
\includegraphics[width=0.95\textwidth]{images/logo9.png}
}


\DeclareRobustCommand{\pag}{%
  \textcolor{pagOrange}{\textbf{\ProjectName}}%
}
%%%%%%%%%%%%%%%%%%%%%%%%%%%%%%%%%%%%%%%%%%%%%%%%%%%%%%%%%%%%



%%%%%%%%%%%%%%%%%%%%%%%%%%%%%%%%%%%%%%%%%%%%%%%%%%%%%%%%%%%%
% Work-package description block. Duplicate the call once per WP.
\newcommand{\WorkPackageDescription}[4]{%
  \par
  \begingroup
  \setlength{\LTpre}{0pt}%
  \setlength{\LTpost}{0pt}%
  %
  \phantomsection
  \addcontentsline{toc}{subsection}{Table 3.1a - #1: #2}%
  \nopagebreak[4]%
  
  \begin{xltabular}{\textwidth}{|
    >{\raggedright\arraybackslash\bfseries}m{0.26\textwidth}|
    X|}
    \hline%
    \cellcolor{pagMidBlue}\color{white}Work package number & #1 \\\hline%
    \cellcolor{pagMidBlue}\color{white}Work package title & #2 \\\hline%
  \end{xltabular}
  \vspace{4pt}
  \begin{xltabular}{\textwidth}{|J{0.99\textwidth}|}
    \hline%
    \patextbf{Objectives.}
    #3\\[1.5em]\hline%
  \end{xltabular}
  \vspace{4pt}\begin{tcolorbox}[
  breakable,
  colback=white,
  colframe=pagBlue,
  boxrule=0.4pt,
  sharp corners,
  boxsep=0pt,
  left=3pt,
  right=3pt,
  top=3pt,
  bottom=3pt,
  before skip=2pt,
  after skip=18pt
]
    \setlength{\emergencystretch}{1em}%
    \patextbf{Description of work.}
    #4
  \end{tcolorbox}

  \endgroup
}
%%%%%%%%%%%%%%%%%%%%%%%%%%%%%%%%%%%%%%%%%%%%%%%%%%%%%%%%%%%%%%%%


\renewcommand{\tabularxcolumn}[1]{>{\raggedright\arraybackslash}p{#1}}
\newcommand{\src}[1]{\vfill{\tiny\color{gray}Source: #1\par}}
\newcommand{\lead}[1]{{\large\color{accent}\textbf{#1}}\par\medskip}
\newcommand{\items}{\setlength{\itemsep}{7pt}}


% Display short DOI values while linking them to their resolver URLs.
\DeclareFieldFormat{doi}{%
  \mkbibacro{DOI}\addcolon\space
  \ifhyperref
    {\href{https://doi.org/#1}{\nolinkurl{#1}}}
    {\nolinkurl{#1}}%
}

\DeclareFieldFormat{year}{\textcolor{DarkOrange}{\bfseries#1}}
\DeclareFieldFormat{date}{\textcolor{DarkOrange}{\bfseries#1}}
\DeclareFieldFormat{labeldate}{\textcolor{DarkOrange}{\bfseries#1}}


\title[Plan4ARI / A4 Review]{Plan4ARI\\\large Review of Activity A4}
\subtitle{Local and Global Autonomous Motion Planning}
\author{Plan4ARI}
\institute{Scientific update and proposed technical specification revision}
\date{Draft for technical discussion / 2 October 2026}
\hypersetup{pdftitle={Plan4ARI~-~Review of Activity A4},pdfauthor={Plan4ARI},pdfsubject={Scientific update and proposed technical specification revision}}

\begin{document}
\begin{frame}
\titlepage
\end{frame}

\begin{frame}{Scope of the review}
\lead{Retained commitments and proposed scientific changes}
\begin{itemize}\items
\item Compare the original A4 methodology and tasks with D4.1 v0.7.
\item Explain the role of MPPI, path-wise redundancy and task localisation.
\item Map the update to deliverables, benchmarks and industrial interfaces.
\item Identify wording and schedule issues for the next specification revision.
\end{itemize}
\medskip
\textbf{Status:} scientific proposals and implementation targets. The repository and design exist; simulator completion, MPPI validation and industrial deployment remain planned.
\src{Technical specification, Activity 4 and project description; D4.1 v0.7, Sections 1, 6, 8 and 10.}
\end{frame}

\begin{frame}{Baseline commitments for A4}
\begin{columns}[T]
\begin{column}{.48\textwidth}
\lead{Application objective}
Autonomous motion planning for teams of manipulators in dynamic, human-shared environments.\par\medskip
Joint and Cartesian planning, with safe velocity adaptation.\par\medskip
MI.RA/Dexter and Core IPC integration.\par
Activity A4: \textbf{M1--M36}.
\end{column}
\begin{column}{.48\textwidth}
\lead{Evaluation commitments}
\begin{itemize}\items
\item \textbf{30\%} reduction in interaction-task execution time.
\item At least \textbf{50 application scenarios}.
\item Mean computation time over \textbf{more than 100 contexts}.
\item OMPL comparison; GPU planners excluded from that comparison.
\end{itemize}
\end{column}
\end{columns}
\src{Technical specification, Activity 4; project description, Expected Results and Computation Time Target.}
\end{frame}

\begin{frame}{Original approach and proposed revision}
\small\renewcommand{\arraystretch}{1.3}
\begin{tabularx}{\textwidth}{@{}p{.20\textwidth}XX@{}}
\toprule
Function & Original approach & D4.1 v0.7 proposal\\\midrule
Global planning & Informed sampling, adaptive local/global exploration, path reuse & Retained; also transfer motions and station feasibility\\
Online adaptation & MPC alongside the sampling planner & MPPI look-ahead over task redundancy in the Open IPC\\
Micro-interpolation & Core-IPC MPC, speed override and sensor interfaces & Deterministic path--velocity MPC with a process speed band\\
Task placement & No explicit pipeline stage & Station/base/rail selection using whole-path feasibility\\
Robot support & Industrial COMAU integration & Parametric robot models, retaining COMAU deployment\\\bottomrule
\end{tabularx}
\src{Technical specification, project methodology and T4.2--T4.5; D4.1, Sections 5--7. Changes proposed for review.}
\end{frame}

\begin{frame}{Scientific evidence behind the update}
\small\renewcommand{\arraystretch}{1.3}
\begin{tabularx}{\textwidth}{@{}p{.25\textwidth}XX@{}}
\toprule
Research development & Implication for A4 & Evaluation required\\\midrule
Vectorised sampling-based planning & Faster CPU kinematics and collision checks for transfers and placement & Same hardware, scenes and budgets as OMPL\\
Tolerance-aware Cartesian planning & Explore tool roll, cone angles and IK classes along a process path & Task accuracy, speed and reconfigurations\\
Constrained and multimodal MPPI & Receding-horizon search with exact IK and separate configuration modes & Constraint violations, latency and fallback tests\\
Predictive path following & Separate geometric path choice from the time law & Speed band and joint-limit satisfaction\\\bottomrule
\end{tabularx}
\src{D4.1, Section 4: VAMP (Thomason et al., 2024), Descartes, constrained MPPI and predictive path-following sources. Literature results are not A4 measurements.}
\end{frame}

\begin{frame}{Scientific contribution to make explicit}
\lead{Joint selection of redundancy, transfers and task placement}
\begin{itemize}\items
\item Explore future tool orientations and external-axis configurations along a task segment.
\item Compare distinct IK solution classes, including the cost of changing class.
\item Plan retract--reconfigure--approach motions and process restarts explicitly.
\item Select stations that make complete segments feasible at the required speed.
\end{itemize}
\medskip
Null-space MPPI alone is established prior work. The proposed contribution concerns its combination with these decisions and industrial execution constraints.
\src{D4.1, Sections 5, 6 and 11; constrained path-integral control and path-wise redundancy literature cited there.}
\end{frame}

\begin{frame}{Reference use case and process constraints}
\begin{columns}[T]
\begin{column}{.48\textwidth}
\lead{Large-structure arc welding}
Industrial arm at a station; mobile base and optional rail extend coverage.\par\medskip
The initial simulator uses a \textbf{fixed COMAU N-220-2.7} and a synthetic stiffened shipbuilding panel.
\end{column}
\begin{column}{.48\textwidth}
\lead{Proposed task parameters}
\begin{itemize}\items
\item Exact seam position; free torch roll.
\item Orientation cone around \textbf{$10^\circ$}.
\item Travel speed band \textbf{$\pm10\%$}.
\item Planned restart overlap around \textbf{1 cm}.
\end{itemize}
\end{column}
\end{columns}
\medskip\small Process experts must confirm tolerances and restart rules. These parameters do not replace the contractual KPIs.
\src{D4.1, Sections 2.3--2.4 and 8.4. Mobile-base configurations are subsequent simulator extensions.}
\end{frame}

\begin{frame}{Proposed functional architecture}
\small\renewcommand{\arraystretch}{1.25}
\begin{tabularx}{\textwidth}{@{}p{.26\textwidth}XX@{}}
\toprule
Stage & Output & Required check\\\midrule
P0 Scene and task model & Registered geometry and constraints & Frames and sensing consistency\\
P1 Task segmentation & Homogeneous process segments & Process continuity and corners\\
P2 Task localisation & Stations, base poses or rail intervals & Whole-segment feasibility\\
P3 Segment planning & IK classes, nominal path and cost-to-go & Collision, limits and speed feasibility\\
P4 Online look-ahead & MPPI path and planned transfers & Mode selection and planning deadline\\
P5 Micro-interpolation & Time law and motion reference & Deterministic limits and safety rules\\\bottomrule
\end{tabularx}
\src{D4.1, Section 6.3. Full research architecture; the initial static mockup implements a smaller scope.}
\end{frame}

\begin{frame}{Multi-goal MPPI and deterministic alternatives}
\lead{Solution classes remain separate throughout optimisation}
\begin{itemize}\items
\item Sample the available task redundancy; compute exact IK within each class.
\item Compare continuation with a class change that includes transfer and restart costs.
\item Use the deterministic segment plan as terminal guidance and a fallback candidate.
\item Report a speed-band violation upstream when the current segment is infeasible.
\end{itemize}
\small
\textbf{Comparison:} deterministic graph/DP replay, predictive IK and null-space NMPC. A fallback requires validation in the current scene; stochastic performance does not establish certification.
\src{D4.1, Sections 6.1--6.2 and 7.}
\end{frame}

\begin{frame}{Open IPC and Core IPC responsibilities}
\small\renewcommand{\arraystretch}{1.3}
\begin{tabularx}{\textwidth}{@{}p{.21\textwidth}XX@{}}
\toprule
Unit & Functions & Interface and timing target\\\midrule
Open IPC & Segmentation, station selection, sampling planner, MPPI and Dexter orchestration & Path $q(s)$, speed band and reconfiguration events; MPPI target 10--100 Hz\\
Core IPC & Path--velocity MPC, hard limits, speed override and control interfaces & Time law $s(t)$, state and constraint events; target 1 kHz\\
A3 sensing & Registration, scene and human updates, seam tracking & Timestamped geometry and uncertainty; synchronous sensor input where required\\\bottomrule
\end{tabularx}\par\medskip
Timing budgets require measurement on the target IPC. The first simulator checks nominal kinematics and does not demonstrate hard-real-time execution.
\src{D4.1, Section 6.4; technical specification, T4.3 and T4.5.}
\end{frame}

\begin{frame}{Proposed revision of T4.1--T4.3}
\small\renewcommand{\arraystretch}{1.3}
\begin{tabularx}{\textwidth}{@{}p{.15\textwidth}XX@{}}
\toprule
Task / period & Proposed clarification & Required evidence\\\midrule
T4.1\newline M1--M6 & Include MPPI, path-wise redundancy, placement and CPU/GPU tool comparison & Scientific decision record and baseline selection\\
T4.2\newline M4--M24 & Retain informed sampling/path reuse; add station checks, transfers and MPPI look-ahead & Ablations, OMPL comparison and constraint checks\\
T4.3\newline M4--M24 & Specify deterministic path--velocity MPC with speed band and upstream infeasibility events & Sensor interfaces, deadline and limit tests\\\bottomrule
\end{tabularx}\par\medskip\small Task periods retained. The detailed informed/local sampling review in D4.1 still requires completion by the T4.2 team.
\src{Technical specification, T4.1--T4.3; D4.1, Sections 4.1, 6, 7 and 10.}
\end{frame}

\begin{frame}{Proposed revision of T4.4--T4.7}
\small\renewcommand{\arraystretch}{1.25}
\begin{tabularx}{\textwidth}{@{}p{.15\textwidth}XX@{}}
\toprule
Task / period & Proposed clarification & Required evidence\\\midrule
T4.4\newline M25--M30 & Wrap P1--P4 as Dockerised Dexter actions with versioned interfaces & Dependency, compatibility and workflow tests\\
T4.5\newline M31--M36 & Integrate P5 in Core IPC; separate asynchronous path input and synchronous feedback & Industrial timing, limit and safety assessment\\
T4.6\newline M4--M36 & Version code, models, scenarios and evaluation configurations & Reproducible releases and regression results\\
T4.7\newline M1--M6 & Initial catalogue/protocol, followed by explicit later evaluation campaigns & Agreed contexts, baselines and operator evaluation\\\bottomrule
\end{tabularx}
\src{Technical specification, T4.4--T4.7; D4.1, Sections 8 and 10. Continued T4.7 evaluation requires schedule alignment.}
\end{frame}

\begin{frame}{Deliverables and schedule dependencies}
\small\renewcommand{\arraystretch}{1.35}
\begin{tabularx}{\textwidth}{@{}p{.15\textwidth}XX@{}}
\toprule
Delivery & Baseline commitment & Review implication\\\midrule
D4.1 / M12 & Scientific plan and first motion-planning mockup & Scientific review at M1--M6 is distinct from mockup delivery; state implemented functions explicitly\\
D4.2 / M24 & Motion-planner module prototypes & Evidence for T4.2/T4.3 prototypes; distinguish research environment from later industrial integration\\
D4.3 / M36 & Motion-planner prototypes integrated in COMAU IPC & Trace Dexter integration at M25--M30 and Core IPC integration at M31--M36\\\bottomrule
\end{tabularx}\par\medskip\small December 2026 and March 2027 are calendar implementation targets. Their correspondence to project months requires the agreed project calendar.
\src{Technical specification, deliverable list and T4.1--T4.5; D4.1, purpose and Section 10.2.}
\end{frame}

\begin{frame}{First mockup and evidence package}
\lead{Initial static-scene planning and nominal trajectory replay}
\begin{itemize}\items
\item Load the robot, cell geometry, seam paths and configured process constraints.
\item Generate a complete plan with baseline adapters before motion starts.
\item Validate collision-free feasibility, Cartesian tolerances, speed and joint limits.
\item Keep the robot stationary and welding inactive if planning fails or times out.
\end{itemize}
\small
\textbf{Review evidence:} software versions, models, seeds, launch procedure, attempt logs, demonstration recording and pending interfaces. Early mockups explicitly identify MPPI as pending.
\src{D4.1, Sections 6.6, 8.4 and 10.2. Planned scope; no completed demonstration is asserted.}
\end{frame}

\begin{frame}{Initial simulation environment design}
\small\renewcommand{\arraystretch}{1.25}
\begin{tabularx}{\textwidth}{@{}p{.25\textwidth}XX@{}}
\toprule
Layer & Proposed implementation & Validation needed\\\midrule
Integration & ROS2, behaviour tree, interchangeable planner adapters & Task sequencing and consistent adapter contracts\\
Scene / replay & MuJoCo; fixed N-220-2.7 and synthetic welding panel & USD conversion, frames, kinematics and collisions\\
Baseline planning & OMPL transfers; Tesseract/Descartes process paths & Compatibility and trajectory validation; later TrajOpt refinement\\
Target platform & Ubuntu 26.04 / ROS2 Lyrical; WSL or native; CPU, up to 32 GB RAM & Dependency compatibility and measured resource use\\
Results & Local dashboard; code/assets in GitHub, campaign data in SharePoint & Reproduction and access to approved assets\\\bottomrule
\end{tabularx}
\src{D4.1, Sections 8.3--8.4; repository design reference dea82b3. Stack choices remain subject to prototype validation.}
\end{frame}

\begin{frame}{Application catalogue and validation coverage}
\small\renewcommand{\arraystretch}{1.25}
\begin{tabularx}{\textwidth}{@{}p{.25\textwidth}Xp{.22\textwidth}@{}}
\toprule
Family & Examples from the 17-application catalogue & Planned evaluation\\\midrule
Welding & Large structures; multi-robot welding & Large-structure welding: hardware + simulation; multi-robot: simulation\\
Continuous toolpaths & Painting, spray, dispensing, WAAM, fibre placement, laser processing & Simulation\\
Surface/contact tasks & Grinding, milling, polishing, NDT & Initial kinematic checks; force/dynamics require extensions\\
Discrete operations & Drilling, riveting, pick-and-place, tending, assembly & Simulation\\
Shared spaces / fleets & Collaborative operations; intralogistics and AMR fleets & Subsequent dynamic-scene evaluation\\\bottomrule
\end{tabularx}\par\medskip\small The December environment covers the first static welding application. It does not complete all 17 families or the 50-scenario commitment.
\src{D4.1, Sections 8.1 and 10.2. Application families summarise the catalogue without changing its scope.}
\end{frame}

\begin{frame}{Baseline KPIs and additional metrics}
\begin{columns}[T]
\begin{column}{.48\textwidth}
\lead{Baseline commitments}
\begin{itemize}\items
\item Interaction execution time: $-30\%$.
\item 50 or more application scenarios.
\item Mean planning time over $>100$ contexts; OMPL baseline.
\item Registration accuracy: about $+20\%$.\item Tracking performance: about $+50\%$.
\end{itemize}
\end{column}
\begin{column}{.48\textwidth}
\lead{Additional D4.1 metrics}
\begin{itemize}\items
\item Success, timeout and failure rates.
\item Planning/replanning latency.
\item Stations, restarts and IK changes.
\item Speed-band and task violations.
\item Jerk, joint margins and CPU/RAM.
\end{itemize}
\end{column}
\end{columns}
\src{Technical specification, A4 expected results; D4.1, Section 8.2. Registration and tracking targets require dedicated sensing/control evidence; static kinematics alone cannot substantiate them.}
\end{frame}

\begin{frame}{Proposed evaluation protocol}
\begin{enumerate}\items
\item Agree on scenario instances, difficulty, constraints and completion criteria.
\item Use identical contexts, hardware budgets and seeds for comparable planners.
\item Report every attempt, including invalid plans, timeouts and failures.
\item Analyse nominal trajectory quality on successful runs, separated by scenario type.
\item Compare MPPI/no MPPI, optimised/fixed placement, horizon length and IK-class count.
\end{enumerate}
\small
Execution-time reduction: $100\,(T_{\mathrm{baseline}}-T_{\mathrm{A4}})/T_{\mathrm{baseline}}$. Report GPU experiments separately; human-interaction metrics follow dynamic-scene and hardware validation.
\src{Technical specification, Computation Time Target; D4.1, Sections 8.2--8.4.}
\end{frame}

\begin{frame}{Implementation targets and research extensions}
\small\renewcommand{\arraystretch}{1.3}
\begin{tabularx}{\textwidth}{@{}p{.22\textwidth}XX@{}}
\toprule
Calendar target & Planned output & Evidence\\\midrule
October 2026 & Validate stack, robot model and synthetic cell & Reproducible loading, frames and baseline paths\\
November 2026 & BT execution, adapters and independent validation & Static task, failure behaviour and initial logs\\
December 2026 & Initial simulator and CPU benchmark environment & Baseline campaign, dashboard and archived results\\
January--February 2027 & MPPI development and integration preparation & Interfaces and common comparison protocol\\
Around March 2027 & Initial MPPI integration/evaluation & Comparative campaign on established scenarios\\\bottomrule
\end{tabularx}
\src{D4.1, Section 10.2. Planning targets, separate from contractual project months. Mobile bases and broader applications follow.}
\end{frame}

\begin{frame}{Interfaces and responsibilities to agree}
\small\renewcommand{\arraystretch}{1.3}
\begin{tabularx}{\textwidth}{@{}p{.25\textwidth}X@{}}
\toprule
Interface & Required agreement\\\midrule
CNR / planning team & Planner design, IK/task descriptions, benchmark adapters and scientific comparisons\\
COMAU / Dexter & Robot assets, process cases, deployment contracts and industrial requirements\\
COMAU / Core IPC & Path streaming, sensor input, timing budgets, fallback and safety assessment\\
A2 / A3 & Task allocation and constraints; scene registration, tracking and uncertainty\\
Control / user evaluation & Actuation and tracking models; operator evaluation and demonstrator acceptance\\\bottomrule
\end{tabularx}\par\medskip\small D4.1 assigns dynamics/control extensions to the later A6 revision, while its architecture links control to A5. Confirm this activity reference before editing the specification.
\src{Technical specification, A4 roles; D4.1, Sections 6, 8.4 and 10.1. Assignments and cross-activity references require confirmation.}
\end{frame}

\begin{frame}{Potential COMAU design contributions}
\small\renewcommand{\arraystretch}{1.3}
\begin{tabularx}{\textwidth}{@{}p{.25\textwidth}XX@{}}
\toprule
Contribution to confirm & Effect on the mockup & Evidence to retrieve\\\midrule
Robot and cell assets & Valid model conversion and realistic collision geometry & Approved assets, versions and model checks\\
Process/use-case definition & Welding paths, tolerances and restart rules & Annotated seams and expert review\\
Industrial interface design & Dexter adapters, IPC allocation and computing budget & Interface specifications and review decisions\\
Demo and acceptance review & Relevant scenarios and measurable completion criteria & Feedback, test records and resulting changes\\\bottomrule
\end{tabularx}\par\medskip\small Existing platform availability is distinct from project-specific contributions. This slide does not certify completed activities or assign effort.
\src{D4.1, Section 9; technical specification, Host Institution role and T4.3--T4.5.}
\end{frame}

\begin{frame}{Specification corrections for discussion}
\small\renewcommand{\arraystretch}{1.25}
\begin{tabularx}{\textwidth}{@{}p{.22\textwidth}XX@{}}
\toprule
Issue & Current evidence & Proposed correction\\\midrule
Output text & A4 paragraph describes natural language to part program & Replace with motion-planning outputs and interfaces\\
Planning layers & MPC described alongside sampling and inside Core IPC & Distinguish Open-IPC MPPI from deterministic P5 MPC\\
T4.4 wording & Motion-planning task refers to deployment of the ``task planning module'' & Correct module reference and define P1--P4 interfaces\\
T4.7 period & M1--M6 includes later platform/end-user evaluation & Separate protocol design from continuing campaigns\\
M48 output & Milestone cites D4.4; deliverable list ends at D4.3/M36 & Confirm intended M48 output before changing the list\\\bottomrule
\end{tabularx}
\src{Technical specification, project methodology, T4.4/T4.7 and milestones; D4.1, Sections 6 and 10. Discussion points only.}
\end{frame}

\begin{frame}{Validation boundaries and open risks}
\begin{itemize}\items
\item Static kinematics checks nominal feasibility; forces, tracking error and stability need dedicated models and tests.
\item The $30\%$ interaction-time KPI requires dynamic scenes and comparable human/safety conditions.
\item Exact IK does not establish complete feasibility: collisions, limits and timing need independent checks.
\item CPU timing targets, mobile registration and mode-switch stability require experimental evidence.
\item System-level safety assessment governs deployment; MPPI costs and simulation alone do not establish certification.
\end{itemize}
\src{Technical specification, A4 methodology and expected results; D4.1, Sections 7, 8.4 and 11.}
\end{frame}

\begin{frame}{Proposed wording for the A4 objective}
\lead{Working text for the next specification revision}
Activity A4 will develop local and global motion planning for robotic teams in dynamic, human-shared environments, integrated into MI.RA/Dexter and the COMAU IPC architecture.\par\medskip
The framework will combine informed sampling and path reuse with task localisation, deterministic segment planning and MPPI look-ahead over task redundancy. A deterministic Core-IPC path--velocity layer will enforce configured execution limits and support sensor-driven adaptation.\par\medskip
Evaluation will retain the contractual scenario, computation-time and interaction-time targets, using reproducible baselines and reporting simulation, dynamic-scene and industrial validation separately.
\src{Proposed wording based on the technical specification, Activity 4, and D4.1 v0.7. Pending discussion; no specification file is changed.}
\end{frame}

\begin{frame}{Decisions required for the revision}
\begin{enumerate}\items
\item Confirm the layered architecture and the scope assigned to T4.2 and T4.3.
\item Agree on mockup acceptance, calendar/project-month alignment and later MPPI evidence.
\item Confirm the static benchmark scope and later dynamic/control activity interfaces.
\item Resolve T4.4 terminology, continuing T4.7 evaluation and the M48 deliverable reference.
\item Assign model/process inputs, industrial integration owners and comparable KPI protocols.
\end{enumerate}
\src{Summary of issues identified through comparison of the current specification and D4.1 v0.7.}
\end{frame}

\begin{frame}{Sources and document status}
\small
\begin{itemize}\items
\item \textbf{Technical specification:} plan4ari\_capitolato.tex and its included project description/plan sections; original A4 tasks, KPIs and deliverables.
\item \textbf{D4.1 v0.7, 2 October 2026:} scientific update, revised architecture, simulation design and implementation targets.
\item \textbf{A2 review presentation:} organisational reference only; A2 claims, methods and submission dates are not transferred to A4.
\item \textbf{Budget v6 tables:} existing WP4 effort values retained in the following reference slides; no resource amendment proposed.
\end{itemize}
\small Literature findings are summarised from D4.1. Candidate software and performance claims require validation during implementation.
\src{Supplied local documents. The technical specification revision will follow after version consolidation.}
\end{frame}

% Existing A4 budget reference slides retained without changing their data.
% Budget analysis: first three of four project years; CNR = AI+RO.
\begin{frame}{Budget analysis: scope and assumptions}


\begin{table}[H]\centering\scriptsize\caption{Mesi-persona per WP~-~intero progetto.}\label{tab:wp-w-tot}
\begin{tabular}{llrrrr}\toprule\textbf{WP} & \textbf{Soggetto} & \textbf{Fund. Research} & \textbf{Ind. Research} & \textbf{Exp. Development} & \textbf{Totale} \\\midrule
WP4 & HI & 3.81 & 62.40 & 0.00 & 66.21 \\
 & AI & 1.96 & 3.91 & 0.00 & 5.87 \\
 & RO & 9.21 & 19.49 & 0.00 & 28.69 \\
 & \textit{Totale WP} & \textit{14.97} & \textit{85.80} & \textit{0.00} & \textit{100.77} \\
\bottomrule\end{tabular}
\end{table}

\end{frame}

\begin{frame}{Budget analysis: Y1}
  
\begin{table}[H]\centering\scriptsize\caption{Mesi-persona per WP~-~anno 1.}\label{tab:wp-w-y1}
\begin{tabular}{llrrrr}\toprule\textbf{WP} & \textbf{Soggetto} & \textbf{Fund. Research} & \textbf{Ind. Research} & \textbf{Exp. Development} & \textbf{Totale} \\\midrule
WP4 & HI & 1.90 & 20.84 & 0.00 & 22.75 \\
 & AI & 0.98 & 0.78 & 0.00 & 1.76 \\
 & RO & 4.60 & 5.85 & 0.00 & 10.45 \\
 & \textit{Totale WP} & \textit{7.48} & \textit{27.47} & \textit{0.00} & \textit{34.95} \\
\bottomrule\end{tabular}
\end{table}
\end{frame}


\end{document}
